\Needspace{20\baselineskip}
\subsection{Maps to a curve}
\label{subsec:curve-map-critical-fibers}

Let $p:X\to C$ be a surjective holomorphic map with connected
fibers, where $X$ is a connected compact complex manifold of
dimension $n\geq1$ and $C$ is a connected compact Riemann surface.
Critical fibers are allowed. Fix any smooth real $(1,1)$-form
$\beta$ on $C$.

\Needspace{14\baselineskip}
\begin{proposition}\label{curve-map:potential-descent}\index{Descent!critical curve fibers}
Pullback gives a bijection
\[
 \PSH(C,\beta)\longrightarrow\PSH(X,p^*\beta).
\]
It preserves the supremum and the ordinary singularity order.
\end{proposition}
\begin{proof}
The critical set is a proper compact analytic subset of $X$.
It has finitely many irreducible components, and $p$ is constant
on each: its differential vanishes on the smooth locus of that
component, and continuity extends the constant to the component.
Thus the set $D\subset C$ of critical values is finite.

Let $U\in\PSH(X,p^*\beta)$. On a base coordinate disk choose a
smooth $a$ with $\ddc a=\beta$. Over $C\setminus D$, the function
$U+a\circ p$ restricts to a plurisubharmonic function, or the
constant $-\infty$, on each compact connected fiber. The maximum
principle makes it constant there. Local holomorphic sections
then give a function $v$ with $U=v\circ p$ and $a+v$ subharmonic.
This is the local argument of \cref{prop:psh-descent-submersion};
it uses compactness of the fibers, not compactness of the
punctured total space.

The function $v$ is bounded above by $\sup_XU$. The removable
singularity theorem extends $a+v$ subharmonically across the
finitely many punctures. The extensions glue, giving
$v\in\PSH(C,\beta)$. Its pullback is $p^*\beta$-plurisubharmonic
by \cref{lem:smooth-disk-criterion}. It agrees with $U$ off the
proper analytic set $p^{-1}(D)$, hence almost everywhere and
everywhere by \cref{lem:qpsh-almost-everywhere-order}.

Conversely pullback preserves the curvature inequality by the
disk criterion. It is not identically $-\infty$, because $p$ is
surjective. Surjectivity also proves uniqueness and preserves
the supremum and every inequality $v\leq w+c$.
\end{proof}

\Needspace{14\baselineskip}
\begin{lemma}\label{curve-map:lelong-pullback}\index{Lelong number!curve map pullback}
Let $f:(\mathbb C^m,0)\to(\mathbb C,0)$ be a nonconstant
holomorphic germ, with $m\geq1$, and let $e=\operatorname{ord}_0f$.
For every quasi-subharmonic germ $v$ at zero,
\[
 \nu(v\circ f,0)=e\nu(v,0).
\]
\end{lemma}
\begin{proof}
A smooth correction of $v$ pulls back to a smooth function, so
we may assume that $v$ is subharmonic. The Taylor expansion of
$f$ gives $|f(z)|\leq M|z|^e$ on a small ball. Choose a unit vector
$\xi$ on which the leading homogeneous term is nonzero. Then
\[
 f(\zeta\xi)=a\zeta^e+O(\zeta^{e+1}),\qquad a\ne0.
\]
For all small $r$, the remainder has modulus at most
$|a|r^e/4$ on $|\zeta|=r$. If $|t|<|a|r^e/4$, Rouch\'e's theorem
shows that $f(\zeta\xi)-t$ has $e$ zeros in $|\zeta|<r$.
Consequently, with $c=|a|/4$,
\[
 \{|t|<cr^e\}\subset f(B(0,r))
                  \subset\{|t|<Mr^e\}.
\]
The suprema of $v$ on these three sets are ordered. Dividing
by $\log r^2<0$, the two disk terms both tend to $e\nu(v,0)$
by the definition in \cref{sec:analytic-types}. The middle term
therefore has the same limit, proving the formula.
\end{proof}

\Needspace{14\baselineskip}
\begin{corollary}\label{curve-map:ideal-types}\index{Multiplier ideal!all-weight curve pullback}
For $v,w\in\QPSH(C)$, the inclusions
$\mathcal I(tv)\subseteq\mathcal I(tw)$ for every $t>0$ hold
if and only if
\[
 \mathcal I(t(v\circ p))\subseteq\mathcal I(t(w\circ p))
 \quad\text{for every }t>0.
\]
In particular pullback preserves and reflects all-weight
multiplier-ideal equivalence.
\end{corollary}
\begin{proof}
On a smooth curve, local modifications are isomorphisms.
By \cref{thm:local-i-inputs}, the base inclusions are therefore
equivalent to $\nu(v,s)\geq\nu(w,s)$ for every $s\in C$.
Let $\pi$ be any local modification of a neighborhood in $X$
with smooth source, and let $y$ be a point of its source.
In a coordinate at $s=p(\pi(y))$, the germ $p\circ\pi-s$ is
nonconstant. Indeed, a modification is an isomorphism on a
dense open set, and a nonconstant holomorphic map cannot be
constant on an open set. Its order $e_y$ is positive and finite.
The preceding lemma gives
\[
 \nu(v\circ p\circ\pi,y)=e_y\nu(v,s),\qquad
 \nu(w\circ p\circ\pi,y)=e_y\nu(w,s).
\]
The base inequalities thus imply all inequalities upstairs,
and the same valuative theorem gives the ideal inclusions.
Conversely the upstairs inclusions give the inequalities for
the identity modification at any $x\in p^{-1}(s)$. Dividing by
$\operatorname{ord}_x(p-s)>0$ recovers the base inequalities.
Apply this argument in both directions for equality.
\end{proof}

\Needspace{14\baselineskip}
\begin{lemma}\label{curve-map:lone-pullback}\index{Compactness!curve map pullback}
If $v_j,v\in\PSH(C,\beta)$ and $v_j\to v$ in $L^1(C)$, then
$v_j\circ p\to v\circ p$ in $L^1(X)$ for every smooth positive
volume form on $X$.
\end{lemma}
\begin{proof}
Put $M_j=\sup_Cv_j$ and $M=\sup_Cv$. Hartogs' lemma applied at
moving maximum points gives $\limsup M_j\leq M$. A subsequence
with limiting suprema below $M$ would give $v$ below that limit
almost everywhere, hence everywhere, a contradiction. Thus
$M_j\to M$.

By surjectivity, $(v_j-M_j)\circ p$ has supremum zero. Normalized
compactness, \cref{thm:normalized-compactness}, makes these
pullbacks precompact in $L^1(X)$. On every relatively compact
product chart over $C\setminus D$, where $D$ is the finite set
of critical values, Fubini's theorem and bounded smooth volume
densities give local $L^1$ convergence to $(v-M)\circ p$.
Every subsequential limit therefore equals this pullback off
$p^{-1}(D)$, and hence almost everywhere, because that analytic
set has smooth volume zero. All subsequential limits agree,
so the entire normalized sequence converges. Adding $M_j\to M$
proves the assertion.
\end{proof}

\Needspace{14\baselineskip}
\begin{lemma}\label{curve-map:analytic-types}\index{Analytic singularity!curve map descent}
A function $v\in\QPSH(C)$ has analytic singularities if and
only if $v\circ p$ has analytic singularities.
\end{lemma}
\begin{proof}
Composing a local analytic presentation with $p$ proves the
forward implication. The pulled-back generators are not all
identically zero, since a nonconstant map to a curve is open.

Conversely fix $s\in C$ and $x\in p^{-1}(s)$. In coordinates,
choose a complex line through $x$ on which $p-s$ has its finite
order $e$ and is nonconstant, as in the Lelong lemma. On that
line write
\[
 (p-s)(\zeta)=h(\zeta)^e,
\]
where $h$ is a local biholomorphism fixing zero; take an analytic
$e$th root of the nonvanishing factor of the germ. The restriction
of $v\circ p$ to the line is not identically $-\infty$, since
$v$ is finite almost everywhere and this map is open. By
\cref{prop:analytic-type-calculus}, it has analytic singularities.
In one variable its presentation is
$r\log|\zeta|^2+O(1)$ with $r\in\mathbb Q_{\geq0}$, obtained by
taking the least order of its finitely many generators.
Changing to $w=h(\zeta)$ gives
\[
 v(s+w^e)=r\log|w|^2+O(1).
\]
Every sufficiently small $t$ has an $e$th root $w$, so
$v(s+t)=(r/e)\log|t|^2+O(1)$. This is an analytic presentation
on the base, including a bounded one when $r=0$.
\end{proof}

\Needspace{18\baselineskip}
\subsection{The curve of fiber components}
\label{subsec:curve-component-factor}

Dropping the connected-fiber hypothesis requires replacing the
base by the curve of fiber components. The following construction
is the curve case of Stein factorization. It also identifies the
multiplicities needed when pulling back curvature measures.

\Needspace{16\baselineskip}
\begin{proposition}\label{curve-map:component-factor}\index{Stein factorization!curve target}\index{Descent!curve of fiber components}
Let $p:X\to C$ be a surjective holomorphic map, where $X$ is a
connected compact complex manifold of dimension $n\geq1$ and
$C$ is a connected compact Riemann surface. There is a
factorization
\[
 X\xrightarrow{q}C'\xrightarrow{r}C,\qquad p=r\circ q,
\]
where $C'$ is a connected compact Riemann surface, $q$ is
surjective and holomorphic with connected fibers, and $r$ is a
finite holomorphic map. The points of $r^{-1}(s)$ correspond to
the connected components of $p^{-1}(s)$, which are the fibers
of $q$ over those points. The factorization is unique up to a
unique isomorphism commuting with $q$ and $r$.

If $t\in r^{-1}(s)$ has ramification index $e_t$, then, as
effective divisors with multiplicities,
\[
 [p^{-1}(s)]=\sum_{t\in r^{-1}(s)}e_t[q^{-1}(t)].
\]
The degree of $r$ is the number of components of a regular
fiber of $p$.
\end{proposition}
\begin{proof}
We first construct the covering over the regular values,
complete it across the critical values, and then extend its
component quotient to $X$.

The critical values form a finite set $D\subset C$, by the
critical-set argument in the proof of
\cref{curve-map:potential-descent}; that argument does not use
connectedness of the fibers. Over $C^\circ=C\setminus D$, the
map $p$ is a proper submersion. The local construction in
\cref{lem:submersion-component-cover} applies there: horizontal
transport over a disk has compact inverse image over its
closure. Compact fibers have finitely many components, and
the resulting local labels are glued by locally constant
permutations. We obtain
\[
 X^\circ\xrightarrow{q^\circ}C'^\circ
             \xrightarrow{r^\circ}C^\circ,
 \qquad X^\circ=p^{-1}(C^\circ),
\]
with $q^\circ$ a holomorphic submersion with connected fibers
and $r^\circ$ a finite unramified covering of degree $d$.
The complement of a proper analytic subset in a connected
complex manifold is connected, so $X^\circ$ and its image
$C'^\circ$ are connected.
Here only the fiber divisor is removed. Locally, the complement
of the zero set of a nonzero holomorphic function in a ball is
path connected: join two outside points through an outside
intermediate point by two complex-line disks. Each restriction
is nonzero because it contains an outside endpoint, and its
isolated zeros can be avoided. Applying this local argument
along a chain of balls proves the required global connectedness.

For each $s\in D$, take a small coordinate disk centered at
$s$ and disjoint from the other critical values. A connected
component of the covering over its puncture is a finite
connected covering of a punctured disk, of some degree $e$.
It has a coordinate $w$ in which $r^\circ(w)=w^e$: the
universal cover of the punctured disk is a half-plane with
projection $z=\exp\zeta$, and the index-$e$ subgroup of its
deck group is generated by $\zeta\mapsto\zeta+2\pi\mathrm i e$.
The coordinate $w=\exp(\zeta/e)$ gives the asserted form.
Attach one point $w=0$ to each such component. There are
finitely many, with their degrees summing to $d$. Gluing these
disks to $C'^\circ$ gives a connected Riemann surface $C'$ and
a finite holomorphic map $r:C'\to C$. The new points have
disjoint disk neighborhoods, so this gluing is Hausdorff.
The inverse image of a small closed base disk is a finite
union of closed disks; over the complement of their interiors,
the original finite covering is proper. Thus $C'$ is compact.

We extend $q^\circ$. Fix $x\in p^{-1}(s)$, $s\in D$. On a
small coordinate ball $B$ about $x$, the image of $p$ lies in
the chosen disk. The complement $B\setminus p^{-1}(s)$ is
connected. Its image under $q^\circ$ therefore lies in a single
punctured disk just attached to $C'$, since those disks are
the connected components of $r^{-1}$ of the punctured base
disk. In its coordinate $w$, the function $w\circ q^\circ$ is
holomorphic and bounded, and
\[
 (w\circ q^\circ)^e=p-s.
\]
The bounded holomorphic removable-singularity theorem extends
it across $p^{-1}(s)\cap B$. To see this without assuming that
the divisor is smooth, first extend at its smooth points by
the one-variable theorem in transverse coordinates. The
singular locus of the reduced divisor has complex codimension
at least two in the ambient manifold, so Hartogs extension
across that locus completes the extension. In dimension one
there are only isolated points. Its $e$th power remains $p-s$ by
the identity theorem, and hence its value is zero on that
fiber. This defines a holomorphic map from $B$ to the attached
disk. The local extensions agree on their dense common
regular locus and therefore glue to $q:X\to C'$ with
$p=r\circ q$. Its image is closed, by compactness, and contains
the dense set $C'^\circ$, so $q$ is surjective.

Every fiber of $q$ is connected. Otherwise write one fiber
as a union of two disjoint nonempty compact subsets and put
them in disjoint open neighborhoods $U_1,U_2$. Properness
gives a disk $V$ about their base point $t$ such that
\[
 q^{-1}(V)\subset U_1\cup U_2.
\]
The map $q$ is open: at any point, a complex line on which
the nonconstant germ $q-t$ has its first nonzero term gives
an open one-variable restriction. Consequently both
$q(U_i)$ contain a neighborhood of $t$. A nearby point of
$C'^\circ$ then has a fiber meeting both neighborhoods,
contradicting connectedness of the regular fibers.

The fibers of $p$ are the disjoint union of these connected
$q$-fibers over the finitely many points of $r^{-1}(s)$.
They are its components. Locally at $t$, the formula for a
finite map is $r(w)-s=w^{e_t}$ after changing the coordinate
by a holomorphic root of a nonvanishing factor. Pulling back
this equality by $q$ gives the divisor multiplicities stated
above. Over $C^\circ$, $r$ is unramified and each sheet is
one component, proving the degree assertion.

Finally, any factorization with the stated properties has
over $C^\circ$ precisely this covering of fiber components.
It is unramified there: a critical point of its finite map,
together with surjectivity of its first factor, would make
$p$ critical over that base point. Thus the component
identification is holomorphic on $C'^\circ$. On every
punctured disk it identifies two completions of the same
finite covering, and the coordinates $w^e$ above extend it
biholomorphically across the added point. It gives a global
isomorphism; density and surjectivity make it unique and
give commutation with both maps.
\end{proof}
